\section{Results}
\label{sec:results}

All results are gaps against the common no-AI baseline, under the shared-spending-share convention unless stated otherwise. Because every reported figure is a ratio of a shocked arm to the baseline, none depends on re-trending except the growth rates in Section~\ref{sec:results-growth}.

\subsection{Output and its composition}
\label{sec:results-output}

Table~\ref{tab:gaps} reports the year-by-year gaps. By \genEndYear{}, GDP is \genShareShiftTfpGdpGapThirtyAbs{} per cent above the no-AI path with the productivity gain and \genShareShiftGdpGapThirtyAbs{} per cent above it without. Both paths build steadily rather than jumping: the productivity arm opens at \genShareShiftTfpGdpGapTwentySixAbs{} per cent in \genStartYear{} and passes \genShareShiftTfpGdpGapTwentyEightAbs{} per cent by 2028, and the automation-only arm opens at \genShareShiftGdpGapTwentySixAbs{} per cent. That gradual profile is the ramp working; with a stepped exponent the whole level effect would arrive in \genStartYear{}.

\begin{table}[htbp]
\centering
\caption{Macroeconomic gaps against the no-AI baseline, per cent}
\label{tab:gaps}
\small
\begin{tabular}{lrrrrr}
\toprule
 & \genStartYear{} & 2027 & 2028 & 2029 & \genEndYear{} \\
\midrule
\multicolumn{6}{l}{\emph{Automation and productivity}} \\
GDP            & \genShareShiftTfpGdpGapTwentySix{} & \genShareShiftTfpGdpGapTwentySeven{} & \genShareShiftTfpGdpGapTwentyEight{} & \genShareShiftTfpGdpGapTwentyNine{} & \genShareShiftTfpGdpGapThirty{} \\
Capital        & \genShareShiftTfpCapitalGapTwentySix{} & \genShareShiftTfpCapitalGapTwentySeven{} & \genShareShiftTfpCapitalGapTwentyEight{} & \genShareShiftTfpCapitalGapTwentyNine{} & \genShareShiftTfpCapitalGapThirty{} \\
Labour input   & \genShareShiftTfpLabourGapTwentySix{} & \genShareShiftTfpLabourGapTwentySeven{} & \genShareShiftTfpLabourGapTwentyEight{} & \genShareShiftTfpLabourGapTwentyNine{} & \genShareShiftTfpLabourGapThirty{} \\
Wage           & \genShareShiftTfpWageGapTwentySix{} & \genShareShiftTfpWageGapTwentySeven{} & \genShareShiftTfpWageGapTwentyEight{} & \genShareShiftTfpWageGapTwentyNine{} & \genShareShiftTfpWageGapThirty{} \\
Consumption    & \genShareShiftTfpConsumptionGapTwentySix{} & \genShareShiftTfpConsumptionGapTwentySeven{} & \genShareShiftTfpConsumptionGapTwentyEight{} & \genShareShiftTfpConsumptionGapTwentyNine{} & \genShareShiftTfpConsumptionGapThirty{} \\
Investment     & \genShareShiftTfpInvestmentGapTwentySix{} & \genShareShiftTfpInvestmentGapTwentySeven{} & \genShareShiftTfpInvestmentGapTwentyEight{} & \genShareShiftTfpInvestmentGapTwentyNine{} & \genShareShiftTfpInvestmentGapThirty{} \\
Tax revenue    & \genShareShiftTfpRevenueGapTwentySix{} & \genShareShiftTfpRevenueGapTwentySeven{} & \genShareShiftTfpRevenueGapTwentyEight{} & \genShareShiftTfpRevenueGapTwentyNine{} & \genShareShiftTfpRevenueGapThirty{} \\
Debt           & \genShareShiftTfpDebtGapTwentySix{} & \genShareShiftTfpDebtGapTwentySeven{} & \genShareShiftTfpDebtGapTwentyEight{} & \genShareShiftTfpDebtGapTwentyNine{} & \genShareShiftTfpDebtGapThirty{} \\
\addlinespace
\multicolumn{6}{l}{\emph{Automation only}} \\
GDP            & \genShareShiftGdpGapTwentySix{} & \genShareShiftGdpGapTwentySeven{} & \genShareShiftGdpGapTwentyEight{} & \genShareShiftGdpGapTwentyNine{} & \genShareShiftGdpGapThirty{} \\
Capital        & \genShareShiftCapitalGapTwentySix{} & \genShareShiftCapitalGapTwentySeven{} & \genShareShiftCapitalGapTwentyEight{} & \genShareShiftCapitalGapTwentyNine{} & \genShareShiftCapitalGapThirty{} \\
Labour input   & \genShareShiftLabourGapTwentySix{} & \genShareShiftLabourGapTwentySeven{} & \genShareShiftLabourGapTwentyEight{} & \genShareShiftLabourGapTwentyNine{} & \genShareShiftLabourGapThirty{} \\
Wage           & \genShareShiftWageGapTwentySix{} & \genShareShiftWageGapTwentySeven{} & \genShareShiftWageGapTwentyEight{} & \genShareShiftWageGapTwentyNine{} & \genShareShiftWageGapThirty{} \\
Consumption    & \genShareShiftConsumptionGapTwentySix{} & \genShareShiftConsumptionGapTwentySeven{} & \genShareShiftConsumptionGapTwentyEight{} & \genShareShiftConsumptionGapTwentyNine{} & \genShareShiftConsumptionGapThirty{} \\
Investment     & \genShareShiftInvestmentGapTwentySix{} & \genShareShiftInvestmentGapTwentySeven{} & \genShareShiftInvestmentGapTwentyEight{} & \genShareShiftInvestmentGapTwentyNine{} & \genShareShiftInvestmentGapThirty{} \\
Tax revenue    & \genShareShiftRevenueGapTwentySix{} & \genShareShiftRevenueGapTwentySeven{} & \genShareShiftRevenueGapTwentyEight{} & \genShareShiftRevenueGapTwentyNine{} & \genShareShiftRevenueGapThirty{} \\
Debt           & \genShareShiftDebtGapTwentySix{} & \genShareShiftDebtGapTwentySeven{} & \genShareShiftDebtGapTwentyEight{} & \genShareShiftDebtGapTwentyNine{} & \genShareShiftDebtGapThirty{} \\
\bottomrule
\end{tabular}
\par\medskip
\begin{minipage}{0.92\textwidth}
\footnotesize \emph{Notes:} per cent difference from the no-AI baseline in each year, shared-spending-share convention. Government spending is a fixed share of output in every arm, so its gap equals the GDP gap by construction and is omitted. Investment is the residual in the model's resource constraint and absorbs numerical error; see Section~\ref{sec:results-output} and Appendix~\ref{app:diagnostics}. Generated from \texttt{results/scenarios.json}.
\end{minipage}
\end{table}

The composition of the gain is the substantive finding. In \genEndYear{} the capital stock is \genShareShiftTfpCapitalGapThirtyAbs{} per cent above baseline in the productivity arm and \genShareShiftCapitalGapThirtyAbs{} per cent above it in the automation-only arm, while aggregate labour input is \genShareShiftTfpLabourGapThirtyAbs{} and \genShareShiftLabourGapThirtyAbs{} per cent \emph{below} it. Both arms raise the return to capital --- from \genBaseInterestThirtyPct{} per cent in the baseline to \genShareShiftTfpInterestThirtyPct{} and \genShareShiftInterestThirtyPct{} per cent respectively --- and cohorts respond by saving more. The extra output is a capital-accumulation result, not a labour-productivity result passed through to workers.

Two features of the automation-only arm deserve attention. First, its GDP ends \genShareShiftGdpGapThirtyAbs{} per cent above baseline even though $Z$ was solved so that output at fixed inputs is unchanged. There is no inconsistency: the neutrality condition in equation~\eqref{eq:zsolve} binds at the baseline's \emph{fixed} $K_0$ and $L_0$, and the higher capital exponent then raises the marginal product of capital, so endogenous saving moves $K$ and output with it. A scenario calibrated to leave the level of GDP unchanged in a model with a fixed capital stock does not leave it unchanged once saving is endogenous, and the \genShareShiftGdpGapThirtyAbs{} per cent gap is the size of that difference in this calibration.

Second, investment should be read with care. It is the residual in the model's resource constraint and absorbs the numerical error reported in Appendix~\ref{app:diagnostics}, so its level gaps --- \genShareShiftTfpInvestmentGapTwentySixAbs{} per cent in \genStartYear{} rising to \genShareShiftTfpInvestmentGapTwentyNineAbs{} per cent in 2029 before falling to \genShareShiftTfpInvestmentGapThirtyAbs{} per cent in \genEndYear{} --- are the least reliable series in the table. The \genEndYear{} fall has a clean interpretation, since the ramp ends that year and the capital stock approaches its new path, so gross investment no longer needs to grow relative to baseline; but the capital stock itself is the series to rely on.

\subsection{Wages, the labour share and consumption}
\label{sec:results-wages}

The labour share follows the imposed path exactly, falling from \genBaseLabourSharePct{} per cent in \genStartYear{} through \genShareShiftTfpLabourShareTwentySeven{}, \genShareShiftTfpLabourShareTwentyEight{} and \genShareShiftTfpLabourShareTwentyNine{} to \genShareShiftTfpLabourShareThirty{} per cent in \genEndYear{}, identically in both arms. This is the mechanical consequence of $\varepsilon = \genEpsilon{}$: the labour share is exactly $1-\gamma_t$, so the model restates the calibration target rather than producing it. The honest reading is that the ramp demonstrates the patch works and that the \genLabourShareFallPp{}-point fall is an assumption carried through the model, not evidence for it.

The real wage falls in both arms --- by \genShareShiftTfpWageGapThirtyAbs{} per cent with the productivity gain and \genShareShiftWageGapThirtyAbs{} per cent without it --- while GDP rises in both. Equation~\eqref{eq:wageidentity} accounts for this exactly. The labour share changes by \genLabourShareLogChangePctAbs{} per cent in logs. Output per worker rises \genShareShiftTfpOutputPerWorkerGapThirtyAbs{} per cent in the productivity arm, which almost but not quite offsets the share effect and leaves the wage \genShareShiftTfpWageGapThirtyAbs{} per cent lower; it rises only \genShareShiftOutputPerWorkerGapThirtyAbs{} per cent in the automation-only arm, leaving the wage \genShareShiftWageGapThirtyAbs{} per cent lower. The result is therefore neither an artefact nor a surprise: in this class of model a labour-share fall of a given size requires an offsetting rise in output per worker of at least the same proportional size before wages can rise, and neither scenario's productivity assumption is large enough to deliver it. In the automation-only arm this is by design, since the OBR's scenario explicitly assumes that higher average labour productivity is \emph{not} reflected in higher real earnings; the model reproduces that feature.

Household consumption, the welfare-relevant aggregate, tracks the wage rather than output. In the productivity arm it is below baseline for most of the window --- \genShareShiftTfpConsumptionGapTwentySixAbs{} per cent in \genStartYear{} --- and ends only \genShareShiftTfpConsumptionGapThirtyAbs{} per cent above it, against GDP \genShareShiftTfpGdpGapThirtyAbs{} per cent above. The gap between the two is the resource cost of the capital accumulation that produces the output gain. In the automation-only arm consumption is below baseline throughout, ending \genShareShiftConsumptionGapThirtyAbs{} per cent lower. Any statement that AI raises UK GDP by \genShareShiftTfpGdpGapThirtyAbs{} per cent in this model is therefore incomplete without the accompanying consumption and wage paths.

\subsection{Growth rates}
\label{sec:results-growth}

Re-trended real GDP growth in \genEndYear{} is \genShareShiftTfpGrowthThirty{} per cent in the productivity arm and \genShareShiftGrowthThirty{} per cent in the automation-only arm, against \genBaseGrowthThirty{} per cent in the baseline. Growth is elevated throughout the ramp rather than in one year, which is the behaviour the time-varying exponent delivers and which a stepped exponent cannot. These are the only results that depend on the trend parameters $g_y$ and $g_n$; every other figure in this section is a ratio and so is invariant to them.

\subsection{The public finances}
\label{sec:results-fiscal}

Under the shared-spending-share convention, tax revenue in \genEndYear{} is \genShareShiftTfpRevenueGapThirtyAbs{} per cent above baseline in the productivity arm, close to but below the GDP gap of \genShareShiftTfpGdpGapThirtyAbs{} per cent, and only \genShareShiftRevenueGapThirtyAbs{} per cent above baseline in the automation-only arm, against a GDP gap of \genShareShiftGdpGapThirtyAbs{} per cent. Revenue grows more slowly than output in both arms, consistent with a shift in the composition of income from labour, which the UK taxes relatively heavily, towards capital. The model reports total tax revenue only, and consumption also falls in the automation-only arm, so the split across tax bases is not identified by these runs. This is the direction of the reduced ``tax-richness'' the OBR describes, though the magnitudes here are small, and the model's effective tax functions are estimated rather than a full statutory calculation, so the gap should be read as directional. Debt is \genShareShiftTfpDebtGapThirtyAbs{} and \genShareShiftDebtGapThirtyAbs{} per cent above baseline in \genEndYear{}: government spending is a fixed share of output in this convention, so a higher output path mechanically raises spending alongside revenue.

Table~\ref{tab:fixed} reports the same scenarios with the level of government spending and transfers held at their baseline paths, the convention the OBR uses when scoring a scenario against a central forecast. The output results are slightly larger --- GDP ends \genFixShareShiftTfpGdpGapThirtyAbs{} and \genFixShareShiftGdpGapThirtyAbs{} per cent above baseline, against \genShareShiftTfpGdpGapThirtyAbs{} and \genShareShiftGdpGapThirtyAbs{} per cent --- because resources not absorbed by higher government spending are available for capital formation. The fiscal result changes, but not uniformly. Debt ends \genFixShareShiftTfpDebtGapThirtyAbs{} per cent \emph{below} baseline in the productivity arm, a reversal of sign against the shared-share run. In the automation-only arm it ends \genFixShareShiftDebtGapThirtyAbs{} per cent above baseline: smaller than the \genShareShiftDebtGapThirtyAbs{} per cent of the shared-share run, but still above it despite spending being held, because revenue rises only \genFixShareShiftRevenueGapThirtyAbs{} per cent. Attributing the remainder would require decomposing the government budget identity, which these runs do not report. Which convention is appropriate depends on the question. The shared-share run holds fiscal policy identical in rule, which is what isolating the AI channel requires, and is the headline here; the fixed-spending run answers the different question of what the scenario does to the public finances if spending plans are held.

\begin{table}[htbp]
\centering
\caption{The same scenarios with government spending and transfers held at baseline levels}
\label{tab:fixed}
\small
\begin{tabular}{lrrrr}
\toprule
 & \multicolumn{2}{c}{Automation and productivity} & \multicolumn{2}{c}{Automation only} \\
\cmidrule(lr){2-3}\cmidrule(lr){4-5}
\genEndYear{} gap, per cent & Shared share & Fixed spending & Shared share & Fixed spending \\
\midrule
GDP          & \genShareShiftTfpGdpGapThirty{} & \genFixShareShiftTfpGdpGapThirty{} & \genShareShiftGdpGapThirty{} & \genFixShareShiftGdpGapThirty{} \\
Capital      & \genShareShiftTfpCapitalGapThirty{} & \genFixShareShiftTfpCapitalGapThirty{} & \genShareShiftCapitalGapThirty{} & \genFixShareShiftCapitalGapThirty{} \\
Wage         & \genShareShiftTfpWageGapThirty{} & \genFixShareShiftTfpWageGapThirty{} & \genShareShiftWageGapThirty{} & \genFixShareShiftWageGapThirty{} \\
Consumption  & \genShareShiftTfpConsumptionGapThirty{} & \genFixShareShiftTfpConsumptionGapThirty{} & \genShareShiftConsumptionGapThirty{} & \genFixShareShiftConsumptionGapThirty{} \\
Tax revenue  & \genShareShiftTfpRevenueGapThirty{} & \genFixShareShiftTfpRevenueGapThirty{} & \genShareShiftRevenueGapThirty{} & \genFixShareShiftRevenueGapThirty{} \\
Debt         & \genShareShiftTfpDebtGapThirty{} & \genFixShareShiftTfpDebtGapThirty{} & \genShareShiftDebtGapThirty{} & \genFixShareShiftDebtGapThirty{} \\
\bottomrule
\end{tabular}
\par\medskip
\begin{minipage}{0.92\textwidth}
\footnotesize \emph{Notes:} per cent difference from the no-AI baseline in \genEndYear{}. Generated from \texttt{results/scenarios.json} and \texttt{results/scenarios\_fixed\_spending.json}.
\end{minipage}
\end{table}

\subsection{Against the quantified scenarios}
\label{sec:results-anthropic}

Table~\ref{tab:anthropic} sets the productivity arm against the ``substantial'' column of Table~3 of \citet{korinek2026}. Two inputs match by construction: the \genLabourShareFallPp{}-point labour-share fall and the \genTfpGainPct{} per cent measured TFP gain. On the outcomes the model undershoots: GDP \genShareShiftTfpGdpGapThirtyAbs{} per cent above the no-AI path against their \genKorinekTableGdpPct{} per cent, growth of \genShareShiftTfpGrowthThirty{} per cent against \genKorinekTableGrowthPct{} per cent, and a capital stock \genShareShiftTfpCapitalGapThirtyAbs{} per cent higher against \genKorinekTableCapitalPct{} per cent. Their unemployment rate of \genKorinekTableUnemploymentPct{} per cent has no counterpart at all.

The undershoot has identifiable sources, and they are structural rather than numerical. \citet{korinek2026} assume an elasticity of substitution of 0.5, against the unit elasticity used here; a lower elasticity makes capital and labour poorer substitutes and changes how much output a given automation intensity delivers. Their framework also carries adoption, capability and task-reinstatement channels that this model collapses into $Z$, and a capital-supply elasticity governing accumulation directly rather than through cohort saving decisions. The comparison therefore bounds how much of their result survives translation into a calibrated fiscal OLG model with a unit elasticity; it does not adjudicate between the two.

\begin{table}[htbp]
\centering
\caption[The productivity arm against Korinek et al. Table 3]{The productivity arm against \protect\citet{korinek2026} Table~3, ``substantial'', \genEndYear{}}
\label{tab:anthropic}
\small
\begin{tabular}{lrrl}
\toprule
Outcome & \citet{korinek2026} (US) & This model (UK) & \\
\midrule
Labour-share fall, pp        & \genKorinekTableLabourShareFallPp{} & \genShareShiftTfpLabourShareFallPp{} & matched by construction \\
Measured TFP, per cent       & \genKorinekTableTfpPct{} & \genTfpGainPct{} & input, matched \\
GDP vs no-AI, per cent       & \genKorinekTableGdpPct{} & \genShareShiftTfpGdpGapThirty{} & model undershoots \\
GDP growth, per cent         & \genKorinekTableGrowthPct{} & \genShareShiftTfpGrowthThirty{} & model undershoots \\
Capital stock, per cent      & \genKorinekTableCapitalPct{} & \genShareShiftTfpCapitalGapThirty{} & model undershoots \\
Unemployment, per cent       & \genKorinekTableUnemploymentPct{} & --- & not modelled \\
\bottomrule
\end{tabular}
\par\medskip
\begin{minipage}{0.92\textwidth}
\footnotesize \emph{Notes:} the comparison is between a US scenario and a UK model, at different elasticities of substitution (0.5 against \genEpsilon{}) and with four of the six AI channels in \citet{korinek2026} absent here. Generated from \texttt{results/scenarios.json} and the Table~3 column recorded in \texttt{uk\_growth\_with\_ai/scenarios.py}.
\end{minipage}
\end{table}

Against the OBR's technological-displacement scenario, three of the four qualitative features are reproduced: the labour share falls, the profit share rises, and revenue grows more slowly than output. The fourth, the rise in equilibrium unemployment to 5.5~per cent, cannot be reproduced in a model without a labour-market block. The OBR's defining quantitative feature --- the level of GDP unchanged against the central forecast --- holds in this model only at fixed inputs, for the reason set out in Section~\ref{sec:results-output}.
